\documentclass[12pt,a4paper]{article}

% --- 核心宏包 ---
\usepackage[utf8]{inputenc}
\usepackage[T1]{fontenc}
\usepackage{ctex} 
\usepackage{amsmath, amsfonts, amssymb}
\usepackage{geometry}
\usepackage{listings}
\usepackage{longtable}
\usepackage{float}
\usepackage{url}
\geometry{left=2.5cm, right=2.5cm, top=3cm, bottom=3cm}

% --- 视觉与排版优化宏包 ---
\usepackage[table]{xcolor} % 颜色支持，带有表格填色功能
\usepackage{booktabs}      % 三线表美化
\usepackage{fancyhdr}      % 自定义页眉页脚
\usepackage[many]{tcolorbox} % 强大的带框/背景盒工具
\usepackage{enumitem}      % 列表间距调整
\usepackage{hyperref}      % 交叉引用与超链接
\usepackage{tikz}
\usetikzlibrary{positioning, arrows.meta} % 必须有 positioning
\definecolor{macblue}{RGB}{0, 122, 255}    % 强调蓝
\definecolor{bglight}{RGB}{245, 245, 247}  % 极简浅灰背景
\definecolor{darkgray}{RGB}{51, 51, 51}
\definecolor{warningred}{RGB}{255, 59, 48} % 增加红色用于攻击/警告的强调
% --- 超链接设置 ---
\hypersetup{
    colorlinks=true,
    linkcolor=macblue,
    urlcolor=macblue
}

% --- 页眉页脚设置 ---
\pagestyle{fancy}
\fancyhf{}
\fancyhead[L]{\textcolor{macblue}{\textbf{Homework}}} % 左侧显示实验报告
\fancyhead[R]{\textbf{RSA 加密算法与解密加速}} % 右侧显示实验名称
\fancyfoot[C]{\thepage}
\renewcommand{\headrulewidth}{0.8pt}
\renewcommand{\headrule}{\hbox to\headwidth{\color{macblue}\leaders\hrule height \headrulewidth\hfill}}

% --- 自定义高亮公式框 ---
\newtcolorbox{coreconcept}[1][]{
    colback=bglight,
    colframe=macblue,
    fonttitle=\bfseries\sffamily,
    title=#1,
    arc=3mm,
    boxrule=1.2pt,
    shadow={2mm}{-1.5mm}{0mm}{black!15},
    left=10pt, right=10pt, top=5pt, bottom=5pt
}
\newtcolorbox{warningbox}[1][]{
    colback=warningred!5,
    colframe=warningred,
    fonttitle=\bfseries\sffamily,
    title=#1,
    arc=3mm,
    boxrule=1.2pt,
    left=10pt, right=10pt, top=5pt, bottom=5pt
}

% --- 标题区域 ---
\title{
\textbf{\LARGE{RSA 加密算法与解密加速}}\\[0.5em]
}
\author{\textbf{202400460046邱泽辉}}
\date{\today}

\begin{document}

\maketitle
\thispagestyle{fancy} % 首页也显示漂亮页眉
\begin{enumerate}
    \item 现有两个 RSA 加密系统，它们的公钥分别为(65537, \(n_1\)) 和(65537, \(n_2\))。
已知模\(n_1 \neq n_2\) 并且\(n_1,n_2\)有大于 1 的公因子。你能在短时间内计算出这两个 RSA
加密系统的私钥吗？请描述计算方法。

解:由于\[
    n_1 = p_1q_1,n_2 = p_2q_2
    \]
其中\(p_1,q_1,p_2,q_2\)都为素数，如果\((n_1,n_2)>1\),那么\((n_1,n_2) = \)\(p_1\)或\(q_1\),否则说明存在因子整除\(p_1\)或者\(q_1\),这与素数矛盾

另一方面，\((n_1,n_2) = \)\(p_2\)或\(q_2\)

不妨设\((n_1,n_2)=p_1=p_2\)

首先使用Eculid算法求出\((n_1,n_2)\),那么\(q_1 = \frac{n_1}{(n_1,n_2)}\)，\(q_2 = \frac{n_2}{(n_1,n_2)}\)

计算\(\phi(n_1)=(p_1-1)(q_1-1),\phi(n_2)=(p_2-1)(q_2-1)\),而\(e=65537\)

使用拓展Eculid算法计算逆元就能得到私钥,\(d_1 = e^{-1} \pmod{\phi(n_1)} \),\(d_2 = e^{-1} \pmod{\phi(n_2)} \)
    \item 对RSA加密算法做如下修改：\(ed\equiv 1 \pmod{\lambda(n)}\)，其中\(\lambda(n)=\frac{(p-1)(q-1)}{gcd((p-1,q-1))}\)
证明修改后方案的正确性。

证明:给定由KeyGen 生成的(pk,sk)，对任意\(m\in Z_n\)，以下等式成立
\[\text{Dec}_{\text{sk}}(\text{Enc}_{\text{pk}}(m))=m\]

由\(ed\equiv 1 \pmod{\lambda(n)}\),存在\(k\)满足\[ed = k\lambda(n)+1\]

则\[
\text{Dec}_{\text{sk}}(\text{Enc}_{\text{sk}}(m)) = \text{Dec}_{\text{sk}}(m^e) = m^{ed} \pmod n
\]

而\[
m^{ed} \pmod n = m^{k\lambda(n)+1}
\]

注意到\[
m^{\lambda(n)} = m^{\text{lcm}(p-1,q-1)}
\]

则
\[
p-1 | \lambda(n)
\]

若\(p\nmid m\),则由费马小定理\(m^{p-1} \equiv 1 \pmod p\),\(m^{\lambda(n)} \equiv 1 \pmod p\)

若\(p| m\)，则\(m^{\lambda(n)}\equiv 1 \pmod p\)

综上，\(m^{\lambda(n)}\equiv 1 \pmod p\)

同理\(m^{\lambda(n)}\equiv 1 \pmod q\)

那么\(m^{\lambda(n)}\equiv 1 \pmod n\)

所以
\(m^{ed} \equiv m \pmod n\)

因此算法正确

\item 已知 n=3026533，利用中国剩余定理加速解密算法，设加密指数 e=3, 解
密指数 d= 2015347，密文 c=152702，求明文 m。

解:因式分解得到\(p=1511,q=2003\),

\(d_p \equiv d \pmod{(p-1)} = 1007\),\(d_q \equiv d \pmod{(q-1)} =1335\)

然后计算
\[
c_p = c \pmod p = 91,c_q = c \pmod q = 474
\]
\[
m_p = c_p^{d_p} \pmod p = 610
\]
\[
m_q = c_q^{d_q} \pmod q = 969
\]
使用中国剩余定理
\[
m = q\times q^{-1} \times m_p + p \times p^{-1} \times m_q \pmod n
\]
其中利用拓展欧几里得算法求逆元
\[
q^{-1} \pmod p = 777,p^{-1} \pmod q = 973
\]
带入得到
\[
m = 2003 \times 777 \times 610 + 1511 \times 973 \times 969 \pmod{3026533} = 1186745
\]
\end{enumerate}
4.\textbf{练习 9.2} \quad 设 $n = pq$（假定 $p > q$）为一个 3076 比特的 RSA 模数。下面的程序表明当 $p - q$ 很小时，此类 RSA 加密可以被破解：
\begin{align*}
&x \leftarrow \lceil 2\sqrt{n} \rceil; \\
&\textbf{while } (\text{true}) \\
\quad &\textbf{if } (x^2 - 4n \text{ is a square of an integer})\ \textbf{break;} \\
\quad &x \leftarrow x + 1; \\
&\textbf{return } (x - 1);
\end{align*}

\begin{enumerate}
\item 已知 $p - q < 10 \cdot 2^{768}$。利用关系
\[
p + q - 2\sqrt{n} = \frac{(p - q)^2}{p + q + 2\sqrt{n}}
\]
（以及 $p + q > 2\sqrt{n}$），证明上面的程序一定在 100 步之内终止。

\item 如果 $S$ 是上面程序的输出，证明
\[
p = \frac{S + \sqrt{S^2 - 4n}}{2}.
\]
\end{enumerate}
解: 
\begin{enumerate}
\item
由于\(p>q,n=pq\),所以
\[
p+q > 2\sqrt{pq} = 2\sqrt{n}
\]
那么一开始\(x\)的初值\(2\sqrt{n}<p+q\)

注意到当\[
x = p+q,x^2-4n = (p+q)^2 - 4n = (p-q)^2
\]
是完全平方数，所以当\(x=p+q\)时程序一定会终止，每一步\(x\)自增1

因此程序运行的步数最多是
\[
p+q - 2\sqrt{n} = \frac{(p-q)^2}{p+q+2\sqrt{n}} < \frac{(p-q)^2}{4\sqrt{n}}
\]
又因为\(n \geq 2^{3075}\)
则
\[
p+q - 2\sqrt{n} < \frac{(p-q)^2}{4\sqrt{n}} < \frac{(10 \cdot 2^{768})^2}{4\cdot 2^{1537.5}} \approx 8.83
\]
所以这个程序在100步内必然可以完成
\item 
不妨设\[
f(x) = x^2 - 4n,x_0 = \lfloor 2\sqrt{n} \rfloor
\]

因为\(n\)不是完全平方数，则\(x_0 < 2\sqrt{n}\)

得先证明\(f(x_0)\)和\(f(p+q)\)之间没有其他完全平方数
\[
f(x_0) < (2\sqrt{n})^2 - 4n = 0 
\]
假设\(f(x)\)为完全平方数，不妨记\(f(x)=y^2(x_0<x<p+q,y>0)\),那么就有
\[
x^2-y^2 = 4n
\]

也就是
\[
(x+y)(x-y) = 4n
\]

由于这个函数是单调递增的，那么
\[
0 < y^2 < f(p+q) = (p-q)^2
\]
那么
\[
0 < y < p-q
\]
\(4n\)只有这么几个因子，\(1,2,4,p,q,2p,2q,4p,4q,n,2n,4n\)

由于\[
2q < x_0 < x+y < 2p
\]

那么只有\[
x+y = n,x-y = 4
\]

那么\[
x = \frac{n+4}{2}
\]

由于\(n+4\)为奇数，那么\(x\)不是整数，矛盾，因此\(f(x_0),f(p+q)\)之间不存在其他完全平方数

因此程序必然在\(x=p+q\)时停止

此时\[
S = p+q,n = pq
\]

\(p,q\)构造二次方程
\[
t^2-St + n = 0
\]
\(p\)是较大的那个根，所以\[
p = \frac{S+\sqrt{S^2-4n}}{2}
\]

\end{enumerate}
\end{document}